%% ============================================================================
%% preamble adalah bagian pengaturan dokumen yang ditulis sebelum \begin{document}.
%% ============================================================================

% ------------------------------------------------------------
% UKURAN KERTAS / MARGIN UNTUK PRINT BOLAK-BALIK
% inner  = sisi jilid
% outer  = sisi luar
% ------------------------------------------------------------
\usepackage[
  a4paper,
  twoside,
  inner=4cm,
  outer=3cm,
  top=3.5cm,
  bottom=3cm,
  bindingoffset=0cm
]{geometry}

% pada main pake settingan ini yaa...
% \documentclass[12pt,a4paper,twoside,openright,fleqn]{report}
% - twoside: Cetak bolak-balik (dua sisi)
% - openright: Bab baru dimulai di halaman ganjil
% - fleqn: opsi LaTeX untuk membuat persamaan matematika rata kiri, bukan berada di tengah halaman.

\usepackage{array}
\usepackage{float}
\usepackage{caption}
\usepackage{enumitem}

% bahasa dokumen
\usepackage[indonesian]{babel}
\usepackage{csquotes}

% butuh ini karena cover pake font Trebuchet MS
\usepackage{fontspec}
% Untuk proposal disertasimu yang membutuhkan font seperti Trebuchet MS dan Times New Roman, 
% saya justru menyarankan memakai XeLaTeX + fontspec, karena bisa menggunakan font sistem secara langsung.

\usepackage{setspace}
\onehalfspacing % Mengaktifkan 1,5 spasi, akan membuat hampir seluruh isi dokumen menggunakan jarak 1,5 spasi.

% --------------------------------------------------
% BAHASA
% --------------------------------------------------
\usepackage[indonesian]{babel}
\usepackage{csquotes}
\usepackage{csquotes}
\DeclareQuoteAlias{english}{indonesian}

% --------------------------------------------------
% SETTING COVER SAMPUL
% --------------------------------------------------
\usepackage{tikz}
\usepackage{graphicx}
\usepackage{xcolor}

% ---------------------------------------------------------
% Font khusus cover
% ---------------------------------------------------------
\newfontfamily\CoverFont{Trebuchet MS}

% --------------------------------------------------
% DAFTAR PUSTAKA
% --------------------------------------------------
\usepackage[
    backend=biber,
    style=authoryear,
    sorting=nyt,
    giveninits=true,
    uniquename=init,
    maxcitenames=2,
    mincitenames=1,
    maxbibnames=99
]{biblatex}

% Karena biblatex belum punya localization Indonesian bawaan
\DeclareLanguageMapping{indonesian}{english}

\addbibresource{isinya/1-umum/referensi.bib}

% --------------------------------------------------
% DAFTAR PUSTAKA - JOURNAL
% --------------------------------------------------

% Format: (Nama, Tahun)
\DeclareDelimFormat{nameyeardelim}{\addcomma\space}

% PENTING: definisikan untuk bahasa INDONESIAN, bukan english
\DefineBibliographyStrings{indonesian}{
    andothers = {dkk\adddot}
}

% Nama daftar pustaka: Neuman, S.P.
\DeclareNameAlias{sortname}{family-given}

% Tidak menampilkan "In:" sebelum jurnal
\renewbibmacro{in:}{}

% Judul artikel dengan tanda kutip
\DeclareFieldFormat[article]{title}{\mkbibquote{#1}}

% Format volume, nomor, halaman
\DeclareFieldFormat[article]{volume}{Vol\adddot\space#1}
\DeclareFieldFormat[article]{number}{No\adddot\space#1}
\DeclareFieldFormat[article]{pages}{hal\adddot\space#1}


% ==================================================
% FORMAT DAFTAR PUSTAKA - BOOK
% ==================================================

% Inisial tanpa spasi: D.K. bukan D. K.
\renewcommand*{\bibinitdelim}{}

% Nama: Todd, D.K.
\DeclareNameAlias{author}{family-given}

% Gunakan "dan", bukan "and"
\DeclareDelimFormat{finalnamedelim}{%
  \ifnumgreater{\value{liststop}}{2}{\finalandcomma}{}%
  \addspace dan\space}

% Edition: 3 -> 3rd
\DeclareFieldFormat{edition}{%
  \ifinteger{#1}
    {\mkbibordinal{#1}}
    {#1}}

% Format buku
\DeclareBibliographyDriver{book}{%
  \usebibmacro{bibindex}%
  \usebibmacro{begentry}%
  \printnames{author}%
  \setunit{\addcomma\space}%
  \printtext[parens]{\printfield{year}}%
  \setunit{\addcomma\space}%
  \printfield[title]{title}%
  \setunit{\addcomma\space}%
  \iffieldundef{volume}
    {}
    {\printfield{volume}%
     \setunit{\addcomma\space}}%
  \iffieldundef{edition}
    {}
    {\printfield{edition}%
     \addspace edition%
     \setunit{\addcomma\space}}%
  \printlist{publisher}%
  \setunit{\addcomma\space}%
  \printlist{location}%
  \usebibmacro{finentry}%
}


% ==================================================
% INPROCEEDINGS
%
% Neuman, S.P. (1980), “Judul”,
% Proceedings ..., Eds: Wang, S.Y. et al.,
% University of Mississippi, Mississippi,
% hal. 189–215.
% ==================================================

\DeclareFieldFormat[inproceedings]{title}{%
  \mkbibquote{#1}}

\DeclareFieldFormat[inproceedings]{booktitle}{%
  \mkbibemph{#1}}

\DeclareFieldFormat[inproceedings]{pages}{#1}

\DeclareNameFormat{proc-editor}{%
  \nameparts{#1}%
  \usebibmacro{name:family-given}
    {\namepartfamily}
    {\namepartgiveni}
    {\namepartprefix}
    {\namepartsuffix}%
  \ifboolexpr{
    test {\ifnumequal{\value{listcount}}{\value{liststop}}}
    and
    test \ifmorenames
  }
    {\addspace et\addabbrvspace al\adddot}
    {}%
}

\DeclareBibliographyDriver{inproceedings}{%
  \usebibmacro{bibindex}%
  \usebibmacro{begentry}%
  \printnames{author}%
  \setunit{\addspace}%
  \printtext[parens]{\printfield{year}}%
  \setunit{\addcomma\space}%
  \printfield{title}%
  \setunit{\addcomma\space}%
  \printfield{booktitle}%
  \setunit{\addcomma\space}%
  \ifnameundef{editor}
    {}
    {%
      \ifboolexpr{
        test {\ifnumgreater{\value{editor}}{1}}
        or
        test {\ifandothers{editor}}
      }
        {\printtext{Eds:\space}}
        {\printtext{Ed:\space}}%
      \printnames[proc-editor]{editor}%
      \setunit{\addcomma\space}%
    }%
  \printlist{organization}%
  \setunit{\addcomma\space}%
  \printlist{location}%
  \setunit{\addcomma\space}%
  \printtext{hal.\addspace}%
  \printfield{pages}%
  \usebibmacro{finentry}%
}
% ==================================================
% FORMAT NAMA EDITOR KHUSUS PROSIDING
% Wang, S.Y. et al.
% ==================================================

\DeclareNameFormat{proc-editor}{%
  \nameparts{#1}%
  \usebibmacro{name:family-given}
    {\namepartfamily}
    {\namepartgiveni}
    {\namepartprefix}
    {\namepartsuffix}%
  \ifboolexpr{
    test {\ifnumequal{\value{listcount}}{\value{liststop}}}
    and
    test \ifmorenames
  }
    {\addspace et\addabbrvspace al\adddot}
    {}%
}

% ==================================================
% GAMBAR
% ==================================================

\usepackage{graphicx}
\usepackage{caption}

% Nama caption
\renewcommand{\figurename}{Gambar}

\graphicspath{{gambar/}} % memberi tahu LaTeX bahwa file gambar secara default dicari di folder ini

\AtBeginDocument{
    % Jarak teks dengan gambar yang berada di tengah halaman
    \setlength{\intextsep}{1\baselineskip}
    
    % Jarak teks dengan gambar yang berada di atas/bawah halaman
    \setlength{\textfloatsep}{0pt}

    % Jarak vertikal antara dua float (gambar/tabel) yang berurutan
    \setlength{\floatsep}{-1\baselineskip}

    % Kalau yang ingin kamu atur adalah jarak setelah caption, gunakan:
    \setlength{\belowcaptionskip}{0\baselineskip}

    % dan kalau ingin mengatur jarak gambar → caption:
    \setlength{\abovecaptionskip}{1\baselineskip}
}

% Format caption gambar
\captionsetup[figure]{
    position=below,
    justification=raggedright,
    singlelinecheck=false,
    font=normalsize,
    labelfont=normalfont,
    textfont=normalfont,
    labelsep=space
}
% ==================================================
% TABEL
% ==================================================
\captionsetup[table]{
    position=top,
    skip=0.7\baselineskip
}


\usepackage{indentfirst} % membuat paragraf pertama setelah judul bab atau subbab ikut menjorok ke dalam.
\setlength\parindent{1.5cm} % mengatur indentasi awal setiap paragraf sebesar 1.5 cm sesuai pedoman
\setlength{\parskip}{0pt} % mengatur tidak ada tambahan jarak vertikal antarparagraf.

% ==================================================
% FONT
% ==================================================
% Jika menggunakan XeLaTeX
\usepackage{fontspec}
\setmainfont{Times New Roman}

% ==================================================
% MATEMATIKA
% ==================================================
% 1. Cek apakah font-nya ada
% Di Terminal jalankan:
% kpsewhich texgyretermes-math.otf

\usepackage{unicode-math}
% \setmathfont{TeX Gyre Termes Math}
\setmathfont{texgyretermes-math.otf}
% Ini lebih cocok daripada Cambria Math karena teks utama kampus Anda menggunakan Times New Roman. 
% TeX Gyre Termes Math secara visual jauh lebih menyatu dengan Times New Roman

\DeclareMathSizes{12}{12}{7}{5}
% Ukuran rumus sesuai pedoman kampus:
% Full = 12 pt
% Subscript/Superscript = 7 pt
% Sub-subscript/Superscript = 5 pt

\usepackage{amsmath}

\numberwithin{equation}{chapter}

% Rumus rata kiri
\setlength{\mathindent}{2pt}

% ==================================================
% JARAK RUMUS
% ==================================================

\makeatletter
\g@addto@macro\normalsize{%
    \setlength{\abovedisplayskip}{0pt}%
    \setlength{\belowdisplayskip}{15pt}%
    \setlength{\abovedisplayshortskip}{0pt}%
    \setlength{\belowdisplayshortskip}{15pt}%
}
\makeatother
\normalsize

% Kemudian Anda cukup menulis:
% \begin{equation}
%     Dice = \frac{2TP}{2TP + FP + FN}
%     \label{eq:dice}
% \end{equation}
% dengan rumus di kiri dan nomor di kanan.
% Untuk merujuknya:
% Persamaan~\eqref{eq:dice}

% ==================================================
% SOURCE CODE / LISTING
% ==================================================
\usepackage{listings}
% Nama caption
\renewcommand{\lstlistingname}{Kode}
% Judul daftar kode
\renewcommand{\lstlistlistingname}{Daftar Kode}

% ==================================================
% HIRARKI PENOMORAN BAB DAN SUBBAB
% ==================================================
\renewcommand{\thesection}{\arabic{chapter}.\arabic{section}}
\renewcommand{\thesubsection}{\thesection.\arabic{subsection}}
\renewcommand{\thesubsubsection}{\thesubsection.\arabic{subsubsection}}
% menentukan sampai mana diberi nomor
\setcounter{secnumdepth}{3}
% menentukan sampai mana muncul di Daftar Isi.
\setcounter{tocdepth}{3}


% ==================================================
% DAFTAR ISI, DAFTAR GAMBAR, DAFTAR TABEL
% ==================================================

\usepackage{tocloft}

% --------------------------------------------------
% NAMA JUDUL
% --------------------------------------------------

\renewcommand{\contentsname}{DAFTAR ISI}
\renewcommand{\listfigurename}{DAFTAR GAMBAR}
\renewcommand{\listtablename}{DAFTAR TABEL}

\usepackage[
    colorlinks=true,
    linkcolor=black,
    citecolor=black,
    urlcolor=blue,
    hidelinks
]{hyperref}

% --------------------------------------------------
% JUDUL DAFTAR ISI
% --------------------------------------------------
\renewcommand{\contentsname}{DAFTAR ISI}

\renewcommand{\cfttoctitlefont}{%
    \fontsize{14pt}{16pt}\selectfont
    \bfseries
}

\renewcommand{\cftaftertoctitle}{}

% Tidak ada tambahan jarak dari margin atas
\setlength{\cftbeforetoctitleskip}{0pt}

% Jarak judul ke isi daftar isi
\setlength{\cftaftertoctitleskip}{1\baselineskip}

\makeatletter
\renewcommand{\@cftmaketoctitle}{%
    \vspace*{\cftbeforetoctitleskip}%
    {%
        \centering
        \cfttoctitlefont
        \contentsname
        \cftaftertoctitle
        \par
    }%
    \vspace{\cftaftertoctitleskip}%
}
\makeatother

% Titik-titik pada BAB
\renewcommand{\cftchapdotsep}{\cftdotsep}
% Jarak antar BAB
\setlength{\cftbeforechapskip}{0.6\baselineskip}
% BAB mulai dari margin kiri
\cftsetindents{chap}{0pt}{1.2em}
% Section lebih masuk
\cftsetindents{section}{1.2em}{2.2em}
% Subsection lebih masuk lagi
\cftsetindents{subsection}{3.6em}{2.5em}
% Subsubsection lebih masuk lagi
\cftsetindents{subsubsection}{7.2em}{4em}

% --------------------------------------------------
% JUDUL DAFTAR GAMBAR
% --------------------------------------------------
\renewcommand{\listfigurename}{DAFTAR GAMBAR}

\renewcommand{\cftloftitlefont}{%
    \fontsize{14pt}{16pt}\selectfont
    \bfseries
}

\renewcommand{\cftafterloftitle}{}

% Tidak ada tambahan jarak dari margin atas
\setlength{\cftbeforeloftitleskip}{0pt}

% Jarak judul ke isi daftar gambar
\setlength{\cftafterloftitleskip}{1\baselineskip}

\makeatletter
\renewcommand{\@cftmakeloftitle}{%
    \vspace*{\cftbeforeloftitleskip}%
    {%
        \centering
        \cftloftitlefont
        \listfigurename
        \cftafterloftitle
        \par
    }%
    \vspace{\cftafterloftitleskip}%
}
\makeatother

% Nomor gambar dimulai dari margin kiri
\setlength{\cftfigindent}{0pt}

% Lebar area nomor gambar (misalnya 1.1, 2.10, dst.)
\setlength{\cftfignumwidth}{2em}

% --------------------------------------------------
% JUDUL DAFTAR TABEL
% --------------------------------------------------

\renewcommand{\listtablename}{DAFTAR TABEL}

\renewcommand{\cftlottitlefont}{%
    \fontsize{14pt}{16pt}\selectfont
    \bfseries
}

\renewcommand{\cftafterlottitle}{}

% Tidak ada tambahan jarak dari margin atas
\setlength{\cftbeforelottitleskip}{0pt}

% Jarak judul ke isi daftar tabel
\setlength{\cftafterlottitleskip}{1\baselineskip}

\makeatletter
\renewcommand{\@cftmakelottitle}{%
    \vspace*{\cftbeforelottitleskip}%
    {%
        \centering
        \cftlottitlefont
        \listtablename
        \cftafterlottitle
        \par
    }%
    \vspace{\cftafterlottitleskip}%
}
\makeatother

% Nomor gambar dimulai dari margin kiri
\setlength{\cfttabindent}{0pt}

% Lebar area nomor gambar (misalnya 1.1, 2.10, dst.)
\setlength{\cfttabnumwidth}{2em}

% ==================================================
% FORMAT BAB, SUB-BAB, SUB-SUB-BAB DAN PARAGRAF
% ==================================================

\usepackage{titlesec}
\usepackage{indentfirst}

% --------------------------------------------------
% PARAGRAF
% --------------------------------------------------

% Awal setiap paragraf masuk 1 tab = 1,5 cm
\setlength{\parindent}{1.5cm}

% Tidak ada jarak tambahan antarparagraf
\setlength{\parskip}{0pt}


% --------------------------------------------------
% CHAPTER / BAB
% Contoh:
%
%                 BAB 1
%
%              PENDAHULUAN
% --------------------------------------------------

\titleformat{\chapter}[display]
    {\centering\fontsize{14pt}{16pt}\selectfont\bfseries}
    {BAB \arabic{chapter}}
    {0\baselineskip}       
    {\MakeUppercase}

\titlespacing{\chapter}
    {0pt}                  % indent kiri
    {0pt}                  % tepat pada batas margin atas
    {2\baselineskip}                  % jarak berikutnya diatur oleh section


% --------------------------------------------------
% SECTION / SUB-BAB
% Contoh:
% 1.1 Latar Belakang
% --------------------------------------------------

\titleformat{\section}[hang]
    {\fontsize{12pt}{14pt}\selectfont\bfseries}
    {\thesection}
    {0.5em}
    {}

\titlespacing{\section}
    {0pt}
    {1.5\baselineskip}       % dari sebelumnya -> judul = 3 spasi / 2x Enter
    {0.5\baselineskip}       % judul -> paragraf = 1,5 spasi / 1x Enter

% --------------------------------------------------
% SUBSECTION
% Contoh:
% 1.1.1 Penyakit Daun Jagung
% --------------------------------------------------

\titleformat{\subsection}[hang]
    {\fontsize{12pt}{14pt}\selectfont\bfseries}
    {\thesubsection}
    {0.5em}
    {}

\titlespacing{\subsection}
    {0pt}
    {1.5\baselineskip}       % sebelum judul = 3 spasi
    {0.5\baselineskip}       % setelah judul = 1,5 spasi


% --------------------------------------------------
% SUBSUBSECTION
% Contoh:
% 1.1.1.1 Northern Corn Leaf Blight
% --------------------------------------------------

\titleformat{\subsubsection}[hang]
    {\fontsize{12pt}{14pt}\selectfont\bfseries}
    {\thesubsubsection}
    {0.5em}
    {}

\titlespacing{\subsubsection}
    {0pt}
    {1\baselineskip}       % sebelum judul = 3 spasi
    {0.3\baselineskip}       % setelah judul = 1,5 spasi


% --------------------------------------------------
% KEDALAMAN PENOMORAN
% --------------------------------------------------

\setcounter{secnumdepth}{3}
\setcounter{tocdepth}{3}